\section{Security Analysis}\label{sec:security}
The contracts are deterministic state machines, so the proofs are short; they establish that the logic meets goals G1--G7 under A1--A5, not that a deployment is secure.

\subsection{Security Properties}\label{sec:props}
\begin{proposition}[Immediate revocation on the audited path, G1]\label{prop:revoke}
If \texttt{revokeDevice(did)} (or \texttt{revokeBatch} containing \texttt{did}) is in a finalized block, every later call to \texttt{evaluateAccess} for \texttt{did} returns DENY.
\end{proposition}
\emph{Proof.} By A1 the call sees finalized state with $\mathrm{status}(\mathrm{did})=$ REVOKED. No function leaves REVOKED, and step~1 of Algorithm~1 runs before every other check and denies all non-ACTIVE statuses. \hfill$\square$

\begin{proposition}[Bounded exposure on the fast path, G1]\label{prop:stale}
Let a revocation finalize at time $t_f$ and let a PEP satisfy A5 (this is an assumption about the deployment, not something the contracts guarantee: if the replica lags more than $\Delta$, the bound fails). A fast-path request evaluated at time $\tau\ge t_f+\Delta$ is denied. If the device submits at most $\rho$ accepted requests per second, at most $\rho\Delta$ requests are accepted after $t_f$.
\end{proposition}
\emph{Proof.} By A5 the replica shows status REVOKED from $t_f+\Delta$ on, and the fast path denies non-ACTIVE devices before any other check. Requests accepted after $t_f$ therefore fall in $[t_f,t_f+\Delta)$, which holds at most $\rho\Delta$ of them. \hfill$\square$

\begin{proposition}[Replay resistance, G2]\label{prop:replay}
For a single PEP and any $(\mathrm{did},c)$ with $c$ a nonce (audited) or counter value (fast path), at most one evaluation returns ALLOW or STEP\_UP.
\end{proposition}
\emph{Proof.} Audited path: the test-and-set of the used-nonce flag is one transaction and transactions execute serially, so a second evaluation sees the flag and denies at step~2. Fast path: the PEP accepts only counters greater than the stored maximum and then stores the new value, so a repeated or older counter is denied. In a multi-PEP zone, the PEP identity must be bound into the signed message so that a request captured at one PEP is useless at another; our testbed exercises a single PEP. \hfill$\square$

\begin{proposition}[Attestation unforgeability, G3]\label{prop:oracle}
A PPT adversary without $sk_\mathcal{O}$ obtains ALLOW or STEP\_UP for a pair $(T,A)$ not attested by $\mathcal{O}$ for that request only with negligible probability.
\end{proposition}
\emph{Proof.} The score checks are reached only after the attestation has not expired and its signature over $(\mathrm{did},\mathrm{res},T,A,t_{\mathrm{valid}},\mathrm{nonce})$ verifies under $pk_\mathcal{O}$. Producing a different $(T,A)$, or a later $t_{\mathrm{valid}}$, with a valid signature is an EUF-CMA forgery (A2). The nonce is inside the message, so an attestation cannot be moved to another request, and the expiry means a good score attested long ago cannot be presented after the device has degraded. Changing only $A$ or only $t_{\mathrm{valid}}$ changes the signed message, which the contract tests demonstrate. \hfill$\square$

\begin{proposition}[Key-reuse prevention, G4]\label{prop:key}
After a finalized revocation, \texttt{registerDevice} with the same key hash reverts.
\end{proposition}
\emph{Proof.} Revocation sets $\mathrm{revokedKey}[\mathrm{H}(pk)]$, and registration reverts on that flag. Since the DID is a function of $\mathrm{H}(pk)$, a new identity needs a new key, which the adversary lacks unless it compromised another device. \hfill$\square$

\begin{proposition}[Governance safety and liveness, G6]\label{prop:gov}
With $N$ voters, quorum $q$, and governance locked: (a) a revocation, reinstatement or batch recall executes only if at least $q$ distinct voters approved it; (b) a set of at most $N-q$ voters cannot stop a proposal that the other $q$ approve; (c) a coalition of at most $q-1$ voters cannot execute any of these actions; (d) no single party, including the administrator, can revoke or reinstate a device, and the most a registrar can do alone is suspend.
\end{proposition}
\emph{Proof.} (a) \texttt{hasVoted} blocks repeats and \texttt{isVoter} blocks outsiders; the approval counter reaches $q$ only with $q$ distinct approvals, and only then is \texttt{revokeDevice} called. (b) Denials end a proposal only when they exceed $N-q$, which at most $N-q$ dissenters cannot, so the $q$ approvers can always complete. (c) Fewer than $q$ approvals never trigger execution. (d) After \texttt{lockGovernance} the registry's governor set is exactly the governance contract, and \texttt{revokeDevice}, \texttt{revokeBatch} and \texttt{reinstateDevice} are the only functions that reach REVOKED or leave SUSPENDED; the contract tests check that the former administrator is rejected. Liveness is conditional: it needs at least $q$ voters to participate. \hfill$\square$

\begin{proposition}[Spoofing, the gate and strictness, G5]\label{prop:bound}
(a) If scores $T,T'$ differ only in signal $s_j$ with weight $w_j$, then $|T-T'|\le w_j$. (b) On a resource with gate $g$, every ALLOW satisfies $A\le g$. (c) On a strict resource, every ALLOW satisfies $T\ge\theta+m$. (d) For an attacker with $P=N=R=1$, $T=1-w_bA$, and it is accepted on a strict resource if and only if $A\le\min\{g,(1-\theta-m)/w_b\}$.
\end{proposition}
\emph{Proof.} (a) $T-T'=w_j(s_j-s_j')$ and $|s_j-s_j'|\le1$. (b) and (c) are steps 5 and 7 of Algorithm~1. (d) Substituting gives $T=0.3+0.3(1-A)+0.2+0.2=1-0.3A$, and ALLOW needs $A\le g$ and $T\ge\theta+m$. \hfill$\square$

\emph{Corollary (what a fixed threshold cannot do).} With the default weights a fully spoofing attacker has $T\ge0.70$ even at $A=1$. Hence on a non-strict resource with $\theta=0.55$ and no gate, \emph{no} anomaly value gets such an attacker denied. For the tiers of Table~\ref{tab:tiers} the bound $(1-\theta-m)/w_b$ is $0.83$ (T2) and $0.60$ (T3), above the gates $0.30$ and $0.20$. The \emph{gate}, not the score, is what stops the attacker, and the residual risk is the set of attackers who keep $A\le g$. The experiments of Section~\ref{sec:gate_eval} measure how that residual behaves.

\subsection{Trust Anchors: Oracle, Voters and Collusion}\label{sec:anchors}
The guarantees rest on two trusted parties that the design does not remove. \emph{The trust oracle} signs the context signals. If its key is stolen, an attacker can attest any posture and network context, which defeats the gate (Proposition~\ref{prop:bound}(d)) but not revocation or replay protection, because the status and nonce checks do not use the oracle. Oracle compromise therefore degrades continuous authorization to signature plus status. A threshold-signed oracle ($t$ of $n$ attesters run by different organizations) with periodic key rotation recorded as a role change is the remedy; we recommend it but have not implemented it. \emph{The voters} are the named consortium members (for example manufacturer, operator and auditor). In our contracts the voter set is fixed at deployment; admitting or removing a voter should itself need the same quorum, which we do not implement. Governance bounds what colluding voters can do: fewer than $q$ cannot revoke any device (Proposition~\ref{prop:gov}), but $q$ or more can revoke any device, which is a denial of service against the plant. Sybil voters are excluded by admission, not by cryptography: voters are identified organizations that join by agreement, so that defence is legal and organizational and the chain cannot enforce it. Choosing $N$ and $q$ is therefore a governance decision, since a larger $N$ raises the collusion threshold and the gas cost (313k--800k gas for $N=5$--$21$).

\subsection{Alignment with the NIST Zero Trust Tenets}\label{sec:tenets}
Table~\ref{tab:tenets} maps the seven tenets of NIST SP~800-207 \cite{nist800207} onto the design and states where it falls short.

\begin{table}[t]
\centering
\caption{Design against the seven NIST SP~800-207 tenets}
\label{tab:tenets}
\scriptsize
\begin{tabular}{@{}p{2.2cm}p{1.1cm}p{3.6cm}@{}}
\toprule
\textbf{Tenet} & \textbf{Status} & \textbf{How, and where it stops} \\
\midrule
1 All data sources and services are resources & Partial & Device-to-resource requests are covered by criticality tier; users and applications are out of scope. \\
2 All communication is secured regardless of location & Not covered & Requests are signed and replay-protected; channel confidentiality (TLS, OPC UA security) is assumed off-chain. \\
3 Access is granted per session & Met, stricter & Decided per request, not per session. \\
4 Access follows dynamic policy (identity, device state, behaviour) & Met & Tiered score of behaviour, posture, network and reliability; posture and network inputs are synthetic in our experiments. \\
5 Integrity and posture of all assets are monitored & Partial & Posture is an oracle-attested input; attestation and monitoring are off-chain dependencies (A3). \\
6 Authentication and authorization are dynamic and enforced before access & Met & Signature, status and policy are checked before each decision; the fast path has the bounded exposure of Proposition~\ref{prop:stale}. \\
7 Information on assets, traffic and access is collected to improve posture & Partial & Decisions are logged and anchored; telemetry collection and policy learning are not evaluated. \\
\bottomrule
\end{tabular}
\end{table}

\subsection{Residual Risks}\label{sec:residual}
The properties say nothing about (i) compromised validators or oracle keys (A1, A4); (ii) a stolen key used by an attacker who behaves like the owner until revoked, which no continuous scheme separates from the owner except by behaviour; (iii) side-channel extraction of keys; or (iv) smart-contract defects, for which the quality-assurance literature recommends testing, verification and audit \cite{wei2024quality}. Also, the trust oracle is a trusted data source, which is the blockchain-oracle problem in a specific form \cite{hassan2023oracle}; a threshold-signed oracle would reduce it (Section~\ref{sec:disc}). All properties hold under A1--A5 (Section~\ref{sec:sysmodel}), and the contracts themselves are not formally verified. Finally, the nonce set costs one storage slot per accepted audited request and needs epoch-based pruning in a long-lived deployment.
